\documentclass[book.tex]{subfiles}
\begin{document}

\chapter{Liquids}
\label{chap:liquids}
\noindent\rule{11cm}{0.4pt} \\
\textbf{Prerequisites:} \autoref{chap:ideal_gas}, \autoref{chap:ising_model}  \\
\textbf{Difficulty Level:} ** \\
\noindent\rule{11cm}{0.4pt}
\vspace{5mm}

Classical fluids are systems of particles which retain a definite volume, and are at sufficiently high temperatures that quantum effects can be neglected. Most liquids at normal temperatures do satisfy this condition,  like liquid air, oil and gasoline, as well as electrolytes, molten salts and salts dissolved in water, which are classified as classical charged fluids.

In this chapter, we develop a simple model to describe the behavior of classical fluids. Unlike ideal gases, whose energy states are purely described by their kinetic energy, the existence of an interactive potential in the case of liquids adds more complexity, as we now have to consider the degree of correlation between fluctuations at different positions. Thus, to determine the distribution function of the correlations from different locations plays a central role in the thermodynamic properties of classical liquids. This can be achieved by x-ray scattering techniques, as we will review at the end of this chapter.

\section{Classical Fluids}

A system of molecules in fluid and solid phases includes all of the particles that comprise the molecules, \textit{i.e.} all protons, neutrons, and electrons. This molecules are often accurately described in terms of the principles of classical mechanics. The behavior of electrons surrounding the nucleus indeed, would most likely be quantum mechanical in nature. However in most cases, electron motion can be averaged over the quantal fluctuations, leading to effective interactions between nuclei. For example, the Lennard-Jones interaction potential is constructed to include two offsetting physical contributions.

\begin{itemize}
\item For short inter-nuclei spacing $r$, the Pauli exclusion principle dictates that the energy diverges exponentially as $r\rightarrow0$. This is approximated in the Lennard-Jones potential as $1/r^{12}$.
\item For large $r$, the electronic polarizability of the neutral atom leads to an attractive interaction that scales as $1/r^6$; such forces are called \textit{Van der Waal's forces}.
\end{itemize}

The Lennard-Jones potential $V_{LJ}(r)$ is given by (Fig.~\ref{fig1}):
\begin{align}\label{Eq1}
V_{LJ}(r)=\nu_0\left[\left(\frac{\sigma}{r}\right)^{12} - \left(\frac{\sigma}{r}\right)^{6}\right]
\end{align}

\noindent where $r$ is the inter-nuclei spacing, $\nu_0$ is a potential scaling term, and $s$ is a length scale of the interactions (related to the polarizability
of the atoms).

Generally, this procedure ignores the coupling between the electron and nucleus kinetic energies; this approximation is valid in the limit $m_e/M\rightarrow 0$ ($m_e \equiv$ electron mass and $M \equiv$ nuclear mass). An important exception where a classical model of a fluid is not acceptable is low temperature helium 

\begin{marginfigure}
\includegraphics[width=\linewidth]{figures/Physics/statistical_mechanics/liquids/Fig1.png}
\caption{Plot of the Lennard-Jones potential $V_{LJ}(r)$ versus the internuclei spacing $r$}
\label{fig1}
\end{marginfigure}

Upon integrating over the electronic motion, we are left with a set of effective particles that interact via an interaction potential, thus defining a classical model for the interactions. 

For an $N$ particle system, the classical degrees of freedom include the coordinate positions 
\begin{align} \{\vec{r}\} = \vec{r}_1, \vec{r}_2,...,\vec{r}_N \end{align} 
\noindent and momenta 
\begin{align} \{\vec{p}\} = \vec{p}_1, \vec{p}_2,...,\vec{p}_N \end{align}
\noindent thus, the microscopic state of the system is defined by the point in phase space $\{ \vec{p}, \vec{r} \}$. For simplicity, we will denote these terms respectively by $\vec{p}, \vec{r}$ when the context is clear.

The thermodynamic behavior is governed by the partition function

\begin{align}
Q = \sum_\mu \exp\left(-\frac{E_\mu}{k_BT}\right) 
\end{align}
\noindent where the summation over $\mu$ implies a sum over quantum states. The energy for a point in phase space is the Hamiltonian $H$, \textit{i.e.}:
\begin{align}\label{Eq2}
E_\mu \rightarrow H(\vec{p}, \vec{r}) = K(\vec{p}) + U(\vec{r})
\end{align}

\noindent where 
\begin{align}
K(\vec{p}) &= \sum_{i=1}^N\frac{\vec{p}_i^2}{2m}
\end{align}
is the total kinetic energy, and $U(\vec{r})$ is the effective classical interaction potential. Points in phase space form a continuum, so the classical partition function must be:
\begin{align}
Q &= A\int d\vec{r}_1d\vec{r}_2...d\vec{r}_N  \int d\vec{p}_1d\vec{p}_2...d\vec{p}_N   \exp\left[-\beta H(\vec{p},\vec{r})\right] 
\end{align}

\noindent where $A$ is a factor to ensure that $Q$ is dimensionless. (Note that we use the alternative style of writing the integral with integration factors at the front rather than at the end for clarity. In general, we will use the mathematical style for integrals with integration factors at the end but may occasionally use the physics style for clarity of expression.)

The factor $A$ corrects for the approximation that the phase space is a continuum, thus $A$ is a universal parameter that does not depend on the Hamiltonian $H$.

We find $A$ by considering the ideal gas Hamiltonian $H = K$, giving the partition function:
\begin{align}
Q &= AV^N\left[\int_{-\infty}^\infty  \int_{-\infty}^\infty \int_{-\infty}^\infty  \exp \left(-\frac{\beta\vec{p}^ 2}{2m}\right) dp_z dp_y dp_x \right]^N\\
&= AV^N\left(\frac{2\pi m}{\beta}\right)^{3N/2} \\ 
&= \frac{1}{N!}\frac{V^N}{\Lambda^{3N}} 
\end{align}

\noindent where the final expression is taken from the translational partition function of ideal gases , and considering that the particles are indistinguishable. 
\begin{align}
\Lambda &= \sqrt{\frac{h^2}{2\pi mk_B T}}
\end{align} 
is the thermal de Broglie wavelength.

Using this equality, we write 
\begin{align}
A &= \frac{1}{N!}\frac{1}{h^{3N}},
\end{align}
and the classical partition function is written as:
\begin{align}
Q &= \frac{1}{N!h^{3N}}\int d\vec{r}_1d\vec{r}_2...d\vec{r}_N \int d\vec{p}_1d\vec{p}_2...d\vec{p}_N \exp\left[-\beta H(\vec{p},\vec{r})\right] \\ 
&= \frac{1}{N!\Lambda^{3N}}\int d\vec{r}_1d\vec{r}_2...d\vec{r}_N \exp[-\beta U(\vec{r})]  \\
&= \frac{1}{N!}\frac{1}{\Lambda^{3N}}Z(T,V,N)  
\end{align}

\noindent where we define the configurational partition function $Z(T,V,N)$, thus leaving only a configurational integral to perform.

The probability of finding the system at locations between $\{\vec{r}\}$ and $\{\vec{r} + d\vec{r}\}$ and momenta between $\{\vec{p}\}$ and $\{\vec{p} + d\vec{p}\}$ system is:
\begin{align}
f(\vec{p},\vec{r})d\vec{r}_1d\vec{r}_2...d\vec{r}_Nd\vec{p}_1d\vec{p}_2...d\vec{p}_N 
\end{align}

\noindent where $f(\vec{p},\vec{r})$ is the probability distribution for observing a system at phase space point $\{\vec{p},\vec{r}\}$.

In this case, the probability distribution has the specific form
\begin{align}
f(\vec{p},\vec{r}) = \frac{1}{Qh^{3N}N!}\exp[-\beta H(\vec{p},\vec{r})] 
\end{align}

The Hamiltonian is 
\begin{align}
H(\vec{p},\vec{r}) = K(\vec{p}) + U(\vec{r}),
\end{align}
thus we can represent the phase space distribution as
\begin{align}
f(\vec{p},\vec{r}) &= \frac{1}{Qh^{3N}N!}e^{-\beta K(\vec{p})}e^{-\beta U(\vec{r})} \\
&= \frac{1}{h^{3N}N!}\frac{e^{-\beta K(\vec{p})}}{\int dp^Ne^{-\beta K(\{ \vec{p}\})}}\frac{e^{\beta U(\vec{r})}}{\int dr^Ne^{\beta U(\vec{r})}} \\
&= \frac{1}{h^{3N}N!}\Phi(p^N)P(r^N)
\end{align}

\noindent where $\Phi(p^N)$ and $P(r^N)$ define the probability distribution for observing system at momentum space point $p^N$ and configuration space point $r^N$, respectively.

We integrate the last expression over positions and momenta of particles 2 to N to get:
\begin{align}
\int d\vec{r}^N \int d\vec{p}_j \prod_{j=2}^N f(\vec{p}, \vec{r})d\vec{p}_1 &= \frac{1}{h^{3N}N!} d\vec{p}_1 \int d\vec{p}_j \prod_{j=2}^N \Phi(p^N) \\ 
&=\frac{1}{h^{3N}N!} \phi(p_1) d\vec{p}_1 \prod_{j=2}^N \int d\vec{p}_j \phi(p_j) \\
&= \frac{1}{h^{3N}N!} \frac{e^{-\beta K(\vec{p}_1)}}{\int d\vec{p}_1 e^{-\beta K( \vec{p}_1)}} d\vec{p}_1 \\
&=(2\pi m k_B T)^{3/2} e^{-\frac{\vec{p}_1^2}{2mk_BT}}d\vec{p}_1
\end{align}

\noindent which is the Maxwell Distribution, that gives the probability that particle 1 has momentum between $\vec{p}_1$ and $\vec{p}_1 + d\vec{p}_1$.

We can rearrange the last expression, which originally considers the differential over the three coordinates of $\vec{p}_1$, into a more compact representation in spherical coordinates, such that $d\vec{p}_1 = dp_{x_1}dp_{y_1}dp_{z_1} = p_1^2dp_1$. Additionally, if we define the expression in terms of the velocity, we obtain a simpler form of the Maxwell Distribution, giving the probability of finding a particle with speed between $v$ and $v + dv$, given by:
\begin{align}\label{Eq3}
f(v)dv = 4\pi m^3 (2\pi mk_BT)^{3/2}\exp\left(-\frac{m}{2k_BT}v^2\right)v^2dv
\end{align}

The most probable speed is $\nu^*= \sqrt{2k_BT/m}$. 

\section{\textit{Equipartition Theorem}}
The equipartition theorem dictates that each thermally active degree of freedom is partitioned a fixed energy.

Assume the Hamiltonian can be diagonalized into $m$ degree of freedom using a general coordinate system $\{q\} = q_1,q_2,...,q_m$, \textit{e.g.} the phonon modes in a crystal, giving the Hamiltonian:

\begin{align}\label{Eq4}
H = \frac{1}{2}\sum_{i=1}^na_ip_i^2 + \frac{1}{2}\sum_{j=1}^mb_jq_j^2
\end{align}

Using the partition function, we can write the average:
\begin{align}
\langle a_i p_i^2\rangle = \frac{\int_{-\infty}^\infty a_ip_i^2e^{-\frac{1}{2}\beta a_ip_i^2}dp_i}{\int_{-\infty}^\infty e^{-\frac{1}{2}\beta a_ip_i^2}dp_i} = \frac{1}{2}k_BT 
\end{align}

\noindent with a similar expression giving $\langle b_jq_j^2\rangle = \frac{1}{2}k_BT$.

Thus, the average energy is $\langle E \rangle = \langle H \rangle = \frac{1}{2}nk_BT + \frac{1}{2}mk_BT$. For a monoatomic gas, $m=0$ and $n=3N$ ($N$ particles and 3 directions), giving $\langle E\rangle = \frac{3}{2}Nk_BT$. For vibration in a solid, $m=3N$ and $n=3N$, giving $\langle E \rangle = 3Nk_BT$.

\section{Pair distribution function}

%\textcolor{red}{TODO: Should this and following section be moved to condensed matter?}

\begin{marginfigure}
\includegraphics[width=\linewidth]{figures/Physics/statistical_mechanics/liquids/Fig2a.png}
\label{fig2a}
\end{marginfigure}
\begin{marginfigure}
\includegraphics[width=\linewidth]{figures/Physics/statistical_mechanics/liquids/Fig2b.png}
\caption{The picture above shows $g(r)$ calculated for a simple simulation of two-dimensional disks. The picture below color codes the particles that comprise the shells surrounding the tagged black particle (see \texttt{http://www.physics.emory.edu/facul ty/weeks//idl/gofr.html} for more details)}
\label{fig2}
\end{marginfigure}

The configurational distribution $P(r^N)$, cannot be factored into single partition functions, as we did with $\Phi(p^N)$ above, because the potential energy $U(\vec{r})$, couples together all of the coordinates. However, we  can analyze the distribution function for a small number of particles by integrating over all coordinates except those corresponding to the particles of interest. For example, define the reduced probability distribution for a pair of particles:
\begin{align}
p^{2,N}(\vec{r}_1,\vec{r}_2) = \frac{1}{Z}\int d\vec{r}_3\int d\vec{r}_4\dotsc \int
d\vec{r}_N\exp[-\beta U(\vec{r})] 
\end{align}

\noindent giving the joint probability of finding particle 1 at $\vec{r}_1$ and particle 2 at $\vec{r}_2$.

Since we don't care about the particle identity, the joint probability 
\begin{align}
\rho^{(2,N)}(\vec{r}_1,\vec{r}_2) &= N(N-1)p^{(2,N)}(\vec{r}_1,\vec{r}_2),
\end{align}
giving the probability for finding any one particle at $\vec{r}_1$ and any other particle at $\vec{r}_2$. Generally, we can write the $n$-particle joint probability as:
\begin{align}
\rho^{(n,N)}(\vec{r}_1,\vec{r}_2,\dotsc,\vec{r}_n) = \frac{N!}{(N-n)!}\frac{1}{Z}\int\prod_{i=n+1}^Nd\vec{r}_i\exp [-\beta U(\vec{r})] 
\end{align}

For an isotropic fluid, 

\begin{align} \rho^{(1,N)}(\vec{r}_1) = \frac{N}{V} = \rho
\end{align}

For a non-interacting system (ideal gas):
\begin{align}
\rho^{(2,N)}(\vec{r}_1,\vec{r}_2) &= \rho^{(1,N)}(\vec{r}_1)\rho^{(1,N)}(\vec{r}_2)\\
&= \frac{N(N-1)}{V^2} \approx \frac{N^2}{V^2} = \rho^2 
\end{align} 

To measure correlations, define the pair distribution function:
\begin{align}\label{Eq5}
g(\vec{r_1},\vec{r_2}) = \frac{\rho^{(2,N)}(\vec{r}_1,\vec{r}_2)}{\rho^2}
\end{align}

The fractional deviation from ideal gas approximation is:
\begin{align}\label{Eq6}
h(\vec{r}_1,\vec{r}_2)=g(\vec{r}_1,\vec{r}_2) - 1
\end{align}

For an isotropic fluid the pair distribution function only depends on $|\vec{r}_1 -\vec{r}_2| = r$, thus $g(\vec{r}_1,\vec{r}_2) = g(r)$ and $ h(\vec{r}_1,\vec{r}_2) = h(r)$.

The pair distribution function gives $\rho g(r)$, the average density of particles at a distance $r$ given that a tagged particle is at the origin (Fig.~\ref{fig2}).

\section{\textit{\Large{Measurement of $g(r)$ by diffraction}}}

%\textcolor{red}{TODO: This should be moved to condensed matter}

Consider a collection of $N$ particles in a fluid. To probe distances $\sim 1\AA$, irradiate the sample with X-rays or neutrons; our discussion will assume X-ray scattering. A beam with wavelength $\lambda$ illuminates the sample. The incident
beam is coherent, \textit{i.e.} all photons are in-phase (Fig.~\ref{fig3}). Due to a difference between the refractive index of the particles and their surroundings, each particle acts as a scattering center

\begin{marginfigure}
% \includegraphics[width=\linewidth]{figures/Physics/statistical_mechanics/liquids/Fig3.png}
\includegraphics[width=\linewidth]{figures/Physics/statistical_mechanics/liquids/Light scattering.png}
\caption{Schematic showing light scattering off of particles in a fluid}
\label{fig3}
\end{marginfigure}

Compare the pathlengths for scattering from 2 different particles (Fig.~\ref{fig4}). The position of the $i$th and $j$th particles are located at $\vec{r}_i$ and $\vec{r}_j$, respectively. The difference in the pathlength of the 2 photons in Fig.~\ref{fig4} is given by the difference between lengths $B$ and $A$, since these are the only segments that differ between the two photons.

Define the incident wavevector:
\begin{align}\label{Eq7}
\vec{q}_0 = \frac{2\pi n}{\lambda}\vec{u}_0
\end{align}

\noindent where $n$ is the index of refraction of the solution, and $\vec{u}_0$ is a unit vector that defines the direction of the incident light. Similarly, define
the scattered wavevector:

\begin{align}\label{Eq8}
\vec{q}_s = \frac{2\pi n}{\lambda}\vec{u}_s
\end{align}

The fixed detector orientation relative to the source is defined by the angle $\theta$.

Using these definitions, the two lengths $B$ and $A$ are given:
\begin{align}\label{Eq9}
B = (\vec{r}_i - \vec{r}_j)\cdot \vec{u}_s
\end{align}

\begin{marginfigure}
% \includegraphics[width=\linewidth]{figures/Physics/statistical_mechanics/liquids/Fig4.png}
% \includegraphics[width=\linewidth]{figures/Physics/statistical_mechanics/liquids/pathlengths for photons.png}
\caption{Schematic showing the difference between pathlengths for
photons scattering from different particles located at $\vec{r}_i$ and $\vec{r}_j$.}
\label{fig4}
\end{marginfigure}

\begin{align}\label{Eq10}
A = (\vec{r}_i - \vec{r}_j)\cdot \vec{u}_0
\end{align}

The pathlength difference is given by:
\begin{align}\label{Eq11}
B - A = (\vec{r}_i - \vec{r}_j)\cdot (\vec{u}_s - \vec{u}_0)
\end{align}

This results in a phase difference given by:
\begin{align}\label{Eq12}
\phi_{ij} = \frac{2\pi n}{\lambda}(B - A) = \frac{2\pi n}{\lambda}(\vec{u}_s - \vec{u}_0) \cdot (\vec{r}_i - \vec{r}_j)
\end{align}

Define the scattering vector $\vec{q}$ as:

\begin{align}\label{Eq13}
\vec{q} = \frac{2\pi n}{\lambda}(\vec{u}_s - \vec{u}_0)
\end{align}

\noindent whose magnitude is:
\begin{align}\label{Eq14}
q = |\vec{q}| = \frac{4\pi n}{\lambda}\sin\left(\frac{\theta}{2}\right)
\end{align}

We have:
\begin{align}\label{15}
\phi_{ij} = \vec{q}\cdot (\vec{r}_i - \vec{r}_j)
\end{align}

The electric field of light scattered by the $i$th particle relative to the phase of light scattered by the $k$th particle is:
\begin{align}\label{Eq16}
E_i = E_0C \cos(2\pi\nu t - \phi _{ik})
\end{align}

\noindent where $n$ is the frequency, $E_0$ is the amplitude of the incident electric field, and $C$ incorporates a number of experiment-specific parameters (\textit{e.g.} distance to detector, wavelength, polarizability).
The intensity of the scattered light is proportional to the total electric field squared. The average scattering intensity $I_s$ is given by the intensity averaged over a complete oscillation period and is given by:
\begin{align}\label{Eq17}
I_s &= 2I_0C^2\nu\int_0^{1/\nu}dt\left[ \sum_{i=1}^N\cos(2\pi\nu t - \phi_{ik}) \right]^2 \\
&= 2I_0C^2\nu\int_0^{1/\nu}dt \sum_{i=1}^N \sum_{j=1}^N\cos(2\pi\nu t - \phi_{ik})\cos(2\pi\nu t - \phi_{jk}) \\
&= 2I_0C^2\nu\int_0^{1/\nu}dt \sum_{i=1}^N \sum_{j=1}^N[\cos(4\pi\nu t - \phi_{ik}- \phi_{jk}) + \cos(\phi_{ik}- \phi_{jk})]
\end{align}

\noindent which makes use of the property:

\begin{align}\label{Eq18}
\cos A \cos B = \frac{1}{2}\cos(A + B) + \frac{1}{2}\cos (A- B)
\end{align}

Noting that:
\begin{align}\label{Eq19}
\int_0^{1/\nu}\cos(4\pi\nu t - \phi_{ik} - \phi_{jk}) dt =0
\end{align}

\noindent and,
\begin{align}\label{Eq20}
\phi_{ik} - \phi_{jk} = \phi_{ij}
\end{align}

\noindent we have,
\begin{align}\label{Eq21}
I_s &= I_0C^2\sum_{i=1}^N\sum_{j=1}^N\cos\phi_{ij} \\
 &= I_0C^2\sum_{i=1}^N\sum_{j=1}^N\cos[\vec{q}\cdot(\vec{r}_i - \vec{r}_j)]
\end{align}

Measuring $I_s$ at zero scattering angle ($\theta = 0$) gives the unknown parameters $I_0$ and $A$. Therefore,

\begin{align}\label{Eq22}
\frac{I_s(\vec{q})}{I_s(\vec{q}=\vec{0})} = \frac{1}{N^2}\sum_{i=1}^N\sum_{j=1}^N\cos[\vec{q}\cdot(\vec{r}_i - \vec{r}_j)]
\end{align}

Define the structure factor $S(\vec{k})$ to be:
\begin{align}\label{Eq23}
S(\vec{k}) = \frac{1}{N^2}\sum_{i=1}^N\sum_{j=1}^N \bigg\langle \exp\left[ i\vec{k}\cdot (\vec{r}_i - \vec{r}_j) \right] \bigg\rangle
\end{align}

\noindent we can write:

\begin{align}\label{Eq24}
\frac{I_s(\vec{q})}{I_s(\vec{q}=\vec{0})} = Re[S(\vec{q})]
\end{align}

Rotational invariance gives $S(\vec{k}) = S(-\vec{k})$, thus:

\begin{align}\label{Eq25}
\frac{I_s(\vec{q})}{I_s(\vec{q}=\vec{0})} = S(\vec{q})
\end{align} 

Expand the structure factor into $i = j$ and $i \neq j$. This gives:
\begin{align}
S(\vec{k}) & = \frac{1}{N^2}N + \frac{1}{N^2}N(N-1)\bigg\langle \exp\left[ i\vec{k}\cdot (\vec{r}_1 - \vec{r}_2) \right] \bigg\rangle \\
& = \frac{1}{N} + \frac{1}{N^2}\frac{N(N-1)\int\prod_{i=1}^Nd\vec{r}_ie^{i\vec{k}\cdot (\vec{r}_1 - \vec{r}_2)}e^{-\beta U}}{\int\prod_{i=1}^Nd\vec{r}_ie^{-\beta U}}\\
&= \frac{1}{N} + \frac{1}{N^2} \int d\vec{r}_1d\vec{r}_2\rho^{(2,N)}(\vec{r}_1,\vec{r}_2)e^{i\vec{k}\cdot (\vec{r}_1 - \vec{r}_2)} \\
&= \frac{1}{N} + \frac{1}{N^2} \int d\vec{r}_1d\vec{r}_{12}\rho^2g(r_{12})e^{i\vec{k}\cdot\vec{r}_{12}} \\
&= \frac{1}{N} + \frac{1}{N}\rho \int d\vec{r}g(r)e^{i\vec{k}\cdot\vec{r}} \\
&= N^{-1} + N^{-1}\rho\tilde{g}(\vec{k}) 
\end{align}

where $\tilde{g}(\vec{k})$ is the Fourier transform of $g(\vec{k})$. This demonstrates that the results of a scattering experiment give a direct measure of the pair distribution of the particles in the fluid.

\end{document}